\documentclass[11pt]{article}

\usepackage[a4paper,margin=1in]{geometry}
\usepackage{amsmath,amssymb,amsthm}
\usepackage{mathtools}
\usepackage{booktabs}
\usepackage{array}
\usepackage{hyperref}
\usepackage{xcolor}
\usepackage{listings}
\usepackage{microtype}

\hypersetup{
  colorlinks=true,
  linkcolor=black,
  urlcolor=blue!50!black,
  pdftitle={Rate-Limit Forecast Model},
  pdfauthor={claumon},
}

\lstset{
  basicstyle=\ttfamily\small,
  frame=single,
  xleftmargin=1em,
  framexleftmargin=0.5em,
  framesep=0.5em,
  framerule=0pt,
  backgroundcolor=\color{gray!10},
}

\newcommand{\rhat}{\hat r}
\newcommand{\sighat}{\hat\sigma}
\newcommand{\rpost}{\hat r_{\text{post}}}
\newcommand{\taupostsq}{\tau_{\text{post}}^2}
\newcommand{\tauzerosq}{\tau_0^2}
\newcommand{\rols}{\hat r_{\text{OLS}}}
\newcommand{\seols}{\mathrm{SE}_{\text{OLS}}}
\newcommand{\seolssq}{\mathrm{SE}_{\text{OLS}}^2}
\newcommand{\sigsessionsq}{\sigma_{\text{session}}^2}
\newcommand{\unow}{u_{\text{now}}}
\newcommand{\tnow}{t_{\text{now}}}
\newcommand{\Treset}{T_{\text{reset}}}
\newcommand{\Dtrem}{\Delta t_{\text{rem}}}
\newcommand{\Cthr}{C_{\text{thr}}}
\newcommand{\taurecent}{\tau_{\text{recent}}}
\newcommand{\Tstar}{T^\ast}
\newcommand{\Tstark}{T^\ast_k}

\title{Rate-Limit Forecast Model}
\author{claumon \texttt{internal/forecast}}
\date{Model version \textbf{v1.1} \enspace --- \enspace 2026-05-28}

\begin{document}
\maketitle

This document specifies the math behind \texttt{internal/forecast}. It is
normative: implementations must match the formulas here, and deviations are
bugs unless this document is updated. The model version above matches
\texttt{forecast.ModelVersion} in the Go code; retired model versions are
preserved under \texttt{internal/forecast/archive/v\textit{N}.\textit{M}/},
and \texttt{CHANGELOG.md} summarises each bump.

\section{What we forecast}

For each open rate-limit window (session, weekly, per-model weekly), the
forecaster reads the time series of utilization snapshots and produces:

\begin{enumerate}
  \item A point estimate $F$ of the utilization at reset.
  \item An 80\% \textbf{credible interval} $[F^-, F^+]$ for $F$. (We use the
    looser ``confidence'' wording in user-facing strings only; the posterior
    is Bayesian, see \S\ref{sec:rate}.)
  \item For each threshold $\Cthr$ (e.g.\ 100\%), either a median ETA with an
    80\% CI, or \texttt{nil} if the threshold is unlikely to be reached
    before reset.
  \item A confidence tag in $\{\text{Low}, \text{Medium}, \text{High}\}$.
\end{enumerate}

This procedure is an \emph{empirical Bayes} forecast: the prior on the rate
$r$ and the path-noise variance $\sigma^2$ are both estimated from past data
(\S\ref{sec:rate}, \S\ref{sec:calib}) and then plugged into the posterior
update, rather than being themselves marginalized under a hyperprior. The
conjugate update on $r$ and the Monte Carlo over rate-and-path uncertainty
are genuinely Bayesian; the ``empirical'' qualifier refers to the prior, not
the inference step.

A separate forecast is produced per gauge. The only cross-window sharing is
through the calibration constants (\S\ref{sec:calib}).

\section{The model in one paragraph}

Inside the current window, utilization $u(t)$ grows at an unknown constant
rate $r$ plus Brownian noise (random behavioral pivots). We estimate $r$ from
the recent slope of the snapshots (OLS) and from a Gaussian prior fit on
past sessions, combined by conjugate update. The forecast at reset is the
sum of two random pieces: deterministic-time-times-uncertain-rate, plus
Brownian accumulation over the remaining horizon. Both pieces are Gaussian,
so the forecast is Gaussian and the CI is a $z$-quantile. The ETA has no
closed form, so we Monte Carlo it.

\section{Generative model and assumptions}

For $s \in [0, \Dtrem]$ with $\Dtrem = \Treset - \tnow$,
\[
  u(\tnow + s) = \unow + r\,s + W(s),
  \qquad W(s) \sim \mathcal{N}\!\bigl(0,\, \sigsessionsq\,s\bigr),
\]
with $r$ unknown and $W$ a Brownian motion independent of $r$. Three
assumptions:

\begin{enumerate}
  \item \textbf{Linear conditional mean.} $r$ is constant within a window.
    Violated by abrupt mid-window mode shifts; partially absorbed by the
    recency window in \S\ref{sec:rate}.
  \item \textbf{Brownian path noise.} Variance grows linearly with time.
    Approximately correct if behavioral pivots are mean-zero and not too
    heavy-tailed.
  \item \textbf{Exchangeable past sessions.} Historical sessions are iid
    draws from the same noise process. Required for the calibration in
    \S\ref{sec:calib}.
\end{enumerate}

\paragraph{On using Brownian motion for a monotone quantity.} Real
utilization is monotone non-decreasing in $[0,1]$: tokens accumulate, the
gauge only resets at the window boundary. A Brownian path can drift below
$\unow$, above $1$, or below $0$, none of which can happen physically. We
keep BM anyway because it yields a Gaussian forecast and closed-form
$z$-quantile CIs; the principled alternative is a non-decreasing process
(Gamma process: $u(t+s) - \unow \sim \mathrm{Gamma}(rs/\beta,\beta)$, or
compound Poisson over individual API calls), at the cost of MC for the CI.

\medskip
The simplification has three concrete consequences:

\begin{itemize}
  \item \textbf{Unphysical lower tail.} $P\!\bigl(u(\tnow+s) < \unow\bigr)
    = \Phi(-r\sqrt{s}/\sigma_{\text{session}})$. With representative
    numbers ($r = 0.06\,\mathrm{h}^{-1}$, $\sigsessionsq = 2.5 \times
    10^{-3}$) this is $\sim 2\%$ at $s = 3\,\mathrm{h}$ but $\sim 20\%$ at
    $s = 30\,\mathrm{min}$. The display clip at $[\unow, 1]$
    (\S\ref{sec:calib}) hides this in user-facing CIs, but the moments
    are biased: the point forecast $F = \unow + \rpost\,\Dtrem$ is
    \emph{slightly low} relative to the truncated-Gaussian mean, with the
    bias worst at short horizons where the lower-tail mass is largest.
  \item \textbf{MC ETAs biased late.} First-passage times in \S6 are
    computed on raw BM paths that can dip below threshold and cross
    later. A monotone process would cross sooner on those paths, so the
    reported median ETA is conservatively late. Magnitude is bounded by
    the same lower-tail probability above.
  \item \textbf{Reinterpretation as forecast-error noise.} A defensible
    reading is that $W(s)$ models the deviation of the realized
    trajectory from the deterministic prediction $\unow + rs$, not literal
    token movement. Under that reading $W(s) < 0$ means ``the user slowed
    down relative to recent rate'', not ``tokens went down''. This is
    consistent with how the calibration in \S\ref{sec:calib} fits $\hat
    a$: against squared forecast errors, not squared step-to-step
    increments.
\end{itemize}

\medskip
The Brownian contribution to $\sigma_F^2$ is $\Dtrem\,\sigsessionsq$
(linear), while the rate-uncertainty contribution from v1.1 is
$\Dtrem^2 \cdot \max(\taupostsq, \bar\tau^2)$ (quadratic), so at long
horizons BM is subdominant. If the new \texttt{UnderspreadX} diagnostic
remains miscalibrated after v1.1 telemetry stabilises, replacing BM with a
Gamma or compound-Poisson process is the natural next investigation.

\paragraph{Visualisation note.} The forecast trajectory modal in the web
UI draws the raw MC paths from \S6, which inherit the BM non-monotonicity
and visibly dip. The dips disappear once a monotone process replaces BM.

\section{Notation}

\begin{center}
\begin{tabular}{@{}ll@{}}
\toprule
Symbol & Meaning \\
\midrule
$u(t),\, \unow$              & utilization at time $t$ and at now, $\in [0,1]$ \\
$\tnow,\, \Treset,\, \Dtrem$ & now, reset, remaining horizon (h) \\
$r$                          & true mean growth rate within window ($\mathrm{h}^{-1}$) \\
$\rols,\, \seols$            & OLS estimate and its standard error \\
$\mu_0,\, \tauzerosq$        & Gaussian prior on $r$ from history \\
$\rpost,\, \taupostsq$       & posterior on $r$ after the OLS update \\
$\sigsessionsq$              & Brownian diffusion coefficient ($\mathrm{h}^{-1}$) \\
$F,\, \sigma_F^2$            & point forecast and its variance \\
$z_{0.9}$                    & $\Phi^{-1}(0.9) \approx 1.2816$ \\
$K,\, \Delta s$              & MC trajectories, step size (defaults: 500, 5 min) \\
\bottomrule
\end{tabular}
\end{center}

\section{Estimating the rate}\label{sec:rate}

Combine two sources of information about $r$.

\paragraph{From the current window.} Let $\mathcal{R}$ be the snapshots in
$[\tnow - \taurecent,\, \tnow]$ and $n = \lvert\mathcal{R}\rvert$. For $n
\geq 3$, fit $u_i = \alpha + r\,t_i + \epsilon_i$ by OLS:
\[
  \rols = \frac{S_{tu}}{S_{tt}},
  \qquad \seolssq = \frac{\sighat_\epsilon^2}{S_{tt}},
\]
with $S_{tt} = \sum (t_i - \bar t)^2$, $S_{tu} = \sum (t_i - \bar t)(u_i -
\bar u)$, and $\sighat_\epsilon^2 = \tfrac{1}{n-2} \sum (u_i - \hat\alpha -
\rols\,t_i)^2$. Default $\taurecent = 30$\,min.

Strictly, Brownian residuals are heteroskedastic and serially correlated, so
the iid-OLS variance formula above is an approximation. For the short
recency window the discrepancy is small and the point estimate is unbiased
either way; we accept the approximation and treat $\seolssq$ as the
likelihood precision in the conjugate update below.

\paragraph{From history.} For each past completed session $s$ with final
value $u_s^\ast$ and duration $D_s$, define $\rho_s = u_s^\ast / D_s$. The
prior mean is the sample mean of $\rho_s$. For the prior variance, note that
under the generative model
\[
  \rho_s = r_s + \frac{W_s(D_s)}{D_s},
  \qquad \mathrm{Var}[\rho_s] = \mathrm{Var}[r_s] + \frac{\sigsessionsq}{D_s},
\]
so the raw sample variance of $\rho_s$ overstates $\mathrm{Var}[r_s]$ by the
average path-noise contribution. Subtract it off:
\[
  \mu_0 = \mathrm{mean}_s\,\rho_s,
  \qquad
  \tauzerosq = \max\!\Bigl(\mathrm{var}_s\,\rho_s
    - \sigsessionsq \cdot \mathrm{mean}_s(1/D_s),\ \varepsilon\Bigr),
\]
with $\varepsilon = 10^{-6}$ guarding against negative estimates from small
samples. This uses the most recent $\sigsessionsq$ from \S\ref{sec:calib};
the two fits depend on each other but only loosely, and the daily refit
cycle converges quickly. On the very first fit (before \S\ref{sec:calib} has
run), $\sigsessionsq = 0$ and the correction is a no-op.

Fit at startup, refresh daily.

\paragraph{Combine.} Normal-normal conjugacy gives the posterior $r \mid
\rols \sim \mathcal{N}(\rpost,\, \taupostsq)$:
\[
  \frac{1}{\taupostsq} = \frac{1}{\tauzerosq} + \frac{1}{\seolssq},
  \qquad
  \rpost = \taupostsq
    \left(\frac{\mu_0}{\tauzerosq} + \frac{\rols}{\seolssq}\right).
\]
Limits: prior dominates when $\seols$ is large (early window); data
dominates when $\seols$ is small (later in the window). For $n < 3$ the OLS
step is undefined; use $\rpost = \mu_0$, $\taupostsq = \tauzerosq$.

\section{Forecast distribution at reset}\label{sec:forecast}

The point forecast is
\[
  F = \unow + \rpost\,\Dtrem.
\]
By the law of total variance, conditioning on $r$:
\[
  \sigma_F^2 = \underbrace{\Dtrem^2 \cdot \max(\taupostsq,\, \bar\tau^2)}_{\text{rate uncertainty}}
            + \underbrace{\Dtrem\,\sigsessionsq}_{\text{path noise}},
\]
where $\bar\tau^2$ is the historical-average rate variance recovered in
\S\ref{sec:calib} (the quadratic coefficient of the §5 regression). The floor
$\max(\taupostsq,\, \bar\tau^2)$ is new in v1.1 and corrects a systematic
under-spread when assumption A1 (constant within-window rate) is violated:
the conjugate $\taupostsq$ then accurately estimates uncertainty about the
\emph{recent} slope but understates uncertainty about the
\emph{end-of-session} slope, because the user pivots. The historical $\bar\tau^2$,
fit from past end-of-session errors, is the right scale for the latter. We
take the max rather than always using $\bar\tau^2$ because for short windows
where $\taupostsq$ is large (few snapshots) the conjugate value should win.

\paragraph{Status as a posterior.} The floor is a robustness correction
rather than a Bayesian update: $\max(\taupostsq,\, \bar\tau^2)$ is not the
variance of any coherent posterior over $r$, so $\sigma_F^2$ should be read
as v1.0's predictive variance with $\taupostsq$ replaced by an empirical
lower bound, not as the variance of a posterior predictive. The marginal
forecast is still Gaussian (we draw $r$ from a Gaussian, then add Gaussian
path noise), with two contributions at different scales: rate uncertainty
is quadratic in horizon (dominates at long horizons), path noise is linear
(dominates at short).

\paragraph{80\% CI.} $[F - z_{0.9}\,\sigma_F,\ F + z_{0.9}\,\sigma_F]$,
clipped to $[\unow, 1]$ for display: utilization within a reset window is
monotone non-decreasing, so the Gaussian left tail dipping below $\unow$ is
unphysical and the displayed lower bound is floored at $\unow$. The unclipped
$F$ is retained for ETA computation.

\section{Calibration of $\sigsessionsq$}\label{sec:calib}

\paragraph{What we learn.} One scalar: how much utilization can drift from
the trend per hour of waiting. It is the only quantity history contributes
to the forecast \emph{spread} (the prior $\mu_0,\,\tauzerosq$ contributes to
the \emph{mean}).

\paragraph{Strategy.} Run the forecaster on past sessions, where we already
know the answer, and measure how wrong it was as a function of horizon.
Read $\sigsessionsq$ off the empirical error structure.

\paragraph{Replay.} For each past session $s$ with reset value $u_s^\ast$ at
$T_s$, sample forecast points $t_f \in [t_s + \taurecent,\, T_s -
\Delta_{\min}]$ (default 6 per session, $\Delta_{\min} = 30$\,min). At each
$t_f$, run \S\ref{sec:rate} on snapshots in $[t_f - \taurecent,\, t_f]$ to
obtain $\rpost(t_f)$, and the projection $\hat F(t_f) = u(t_f) +
\rpost(t_f)\,(T_s - t_f)$. Record
\[
  e_f = u_s^\ast - \hat F(t_f),
  \qquad
  \delta_f = T_s - t_f.
\]

\paragraph{Two sources of error in each residual.} The residual $e_f$ mixes
(A) the rate estimate being off at $t_f$, with the error compounded over
$\delta_f$, and (B) Brownian path noise accumulating over $\delta_f$. By the
same decomposition as \S5,
\[
  \mathbb{E}[e_f^2 \mid \delta_f]
  = \underbrace{\delta_f\,\sigsessionsq}_{\text{path noise (B)}}
  + \underbrace{\delta_f^2\,\bar\tau^2}_{\text{rate uncertainty (A)}},
\]
with $\bar\tau^2$ the average posterior rate variance at the forecast
points. The two contributions scale differently in $\delta_f$, which is what
lets us separate them. Naive averaging like $\sigsessionsq \approx
\mathrm{mean}_f(e_f^2 / \delta_f)$ would not separate them and would bias
the estimate upward by a $\delta_f\,\bar\tau^2$ contamination.

\paragraph{Joint regression.} Fit both unknowns by OLS of $e_f^2$ on
$[\delta_f,\, \delta_f^2]$ (no intercept):
\[
  (\hat a, \hat b)
    = \arg\min_{a, b} \sum_f \bigl(e_f^2 - a\,\delta_f - b\,\delta_f^2\bigr)^2,
  \qquad
  \sigsessionsq := \max(\hat a,\, 10^{-6}).
\]
The coefficient $\hat a$ on the linear term is the per-hour path-noise
variance: that is the quantity we want.

Note that this regression treats $\bar\tau^2$ as constant across forecast
points, while in reality $\taupostsq(t_f)$ varies (it depends on how many
recent snapshots existed at $t_f$). If $\taupostsq(t_f)$ correlates with
$\delta_f$ --- and it usually does, since $\delta_f$ is anti-correlated with
elapsed time in the session --- the missing covariance leaks into $\hat a$.
A tighter alternative is to subtract the known $\delta_f^2\,\taupostsq(t_f)$
piece as an offset and regress the residual on $\delta_f$ alone; we keep
the joint regression for simplicity and absorb the bias into the daily
refit.

\paragraph{Role of $\hat b$ (v1.1: kept as a floor).} Conceptually $\hat b$
estimates the same quantity as $\taupostsq$ in \S\ref{sec:rate} (the rate
variance relevant to predicting $u(\Treset)$), just as a historical average
instead of a fresh per-forecast value. The original v1.0 design discarded
$\hat b$ on the grounds that the per-forecast $\taupostsq$ uses today's
actual snapshots and is therefore strictly more informative.

In practice, on real sessions, $\taupostsq$ shrinks to one or two orders of
magnitude below $\hat b$ within a few snapshots of OLS. Empirically, the
$\hat b$-sized variance is the one that matches the realized end-of-session
errors. The reason is assumption A1: $\taupostsq$ measures uncertainty in
the \emph{recent} slope, which OLS can pin down precisely; but if the rate
shifts mid-window, the slope governing the rest of the session has variance
closer to the cross-session $\bar\tau^2$ than to today's $\taupostsq$.

v1.1 therefore keeps $\hat b$ and stores it as $\bar\tau^2$ on the
calibration. \S\ref{sec:forecast} uses $\max(\taupostsq,\, \bar\tau^2)$ in
$\sigma_F^2$. The MC in \S6 likewise samples $r_k$ with standard deviation
$\sqrt{\max(\taupostsq,\, \bar\tau^2)}$.

The $b\,\delta_f^2$ term has to be present in the regression in any case:
without it, the $\delta_f^2\,\bar\tau^2$ piece of $e_f^2$ has nowhere to go
but $\hat a$, biasing it upward.

\paragraph{Floor and refit.} The floor on $\hat a$ guards against negative
estimates from small samples (OLS does not enforce sign). Refit daily.

\section{Threshold ETA}

For threshold $\Cthr$ with $\unow < \Cthr \leq 1$, the first-passage time is
$\Tstar = \inf\{t > \tnow : u(t) \geq \Cthr\}$, with $\Tstar = \infty$ if no
crossing in $[\tnow,\, \Treset]$.

\paragraph{Deterministic ETA (point estimate).} Setting $W \equiv 0$ and $r
= \rpost$:
\[
  \tilde\Tstar = \tnow + \frac{\Cthr - \unow}{\rpost},
\]
defined when $\rpost > 0$ and $\tilde\Tstar \leq \Treset$; else
\texttt{nil}. Note that $\tilde\Tstar$ is \textbf{not} the median of
$\Tstar$ under the full stochastic model. Two effects pull in opposite
directions:

\begin{itemize}
  \item \emph{BM skew (pulls median down).} For BM with fixed drift $\mu > 0$
    hitting $a > 0$, $\Tstar \sim \mathrm{IG}(a/\mu,\, a^2/\sigma^2)$ with
    mean $a/\mu$ and right skew, so $\mathrm{median}(\Tstar \mid r{=}\mu) <
    a/\mu$.
  \item \emph{Rate uncertainty (pulls median up).} Mixing over $r \sim
    \mathcal{N}(\rpost,\, \taupostsq)$ inflates first-passage times
    asymmetrically: the map $r \mapsto a/r$ is convex on $r > 0$ and
    steepens sharply as $r \to 0^+$, so the right tail of $\Tstar$
    stretches far more than the left tail compresses. Any mass at $r \leq
    0$ contributes additional atoms at $\Tstar = \infty$. Both effects pull
    the median of $\Tstar$ above $\tilde\Tstar$.
\end{itemize}

Which effect wins depends on the regime, so $\mathrm{median}(\Tstar)$ can
sit on either side of $\tilde\Tstar$. The MC percentiles below are the
authoritative summary; $\tilde\Tstar$ is a cheap deterministic anchor, not
an estimator of the median.

\paragraph{CI via Monte Carlo.} The first-passage time of a Brownian motion
with random Gaussian drift has no tractable closed form. Simulate $K = 500$
trajectories:

\begin{lstlisting}
tau_eff^2 = max(tau_post^2, bar_tau^2)     # v1.1 floor; see §4
for k in 1..K:
    r_k    = N(r_hat_post, tau_eff^2)      # may be negative; do NOT clip
    u_k[0] = u_now
    for j in 1..M, dt = step size of jth bucket:
        u_k[j] = u_k[j-1] + r_k * dt + N(0, sigma_session^2 * dt)
    T*_k = linear interpolation of first j with u_k[j] >= C_thr
         = infinity if none
\end{lstlisting}

Negative $r_k$ draws are kept as-is: those trajectories drift downward and
typically yield $\Tstark = \infty$, which is the correct outcome under the
model. Clipping to zero would bias crossings upward and contradict \S5.

\paragraph{Reported summaries.} Let $p_\infty$ be the fraction of
trajectories with $\Tstark = \infty$. Define $\tilde\Tstar_{\text{MC}} =
\mathrm{median}(\Tstark)$ over the \textbf{full} sample (treating $\infty$
as larger than any finite value):

\begin{itemize}
  \item If $p_\infty \geq 0.5$: the true median is $\infty$. Report ETA as
    \texttt{nil}.
  \item If $p_\infty < 0.5$ but $p_\infty \geq 0.1$: the median is finite
    but the 90th percentile is $\infty$. Report median, lower bound = 10th
    percentile of finite $\Tstark$, upper bound = \texttt{nil}
    (open-ended).
  \item If $p_\infty < 0.1$: report median and full 80\% CI from the 10th
    and 90th percentiles of finite $\Tstark$.
\end{itemize}

RNG seed derived from $(\tnow,\, \Treset,\, \unow,\, \rpost,\, \taupostsq,\,
\sigsessionsq)$ for test reproducibility.

\section{Confidence tag}

Summarises the effective amount of evidence:
\[
  N_{\text{eff}} = \min\!\left(n,\
    \frac{\tauzerosq}{\seolssq} + \lvert\mathcal{S}\rvert\right).
\]

\begin{center}
\begin{tabular}{@{}ll@{}}
\toprule
Tag & Threshold \\
\midrule
High   & $N_{\text{eff}} \geq 50$ \\
Medium & $15 \leq N_{\text{eff}} < 50$ \\
Low    & $N_{\text{eff}} < 15$ \\
\bottomrule
\end{tabular}
\end{center}

Thresholds are heuristics; revisit after seeing real values.

\section{Edge cases}

\begin{center}
\begin{tabular}{@{}p{0.4\textwidth}p{0.55\textwidth}@{}}
\toprule
Case & Handling \\
\midrule
$n < 3$ snapshots in $\mathcal{R}$
  & Use prior alone: $\rpost = \mu_0$, $\taupostsq = \tauzerosq$ \\
Prior empty ($\lvert\mathcal{S}\rvert < 2$)
  & Suppress forecast; UI shows ``collecting data'' \\
Reset detected in $\mathcal{R}$ ($u_{i+1} < u_i$)
  & Drop pre-reset snapshots; refit \\
$\Cthr \leq \unow$
  & Threshold already crossed; return $\Tstar = \tnow$ \\
$\rpost \leq 0$
  & Deterministic ETA undefined; rely on MC summary \\
$F \notin [0, 1]$
  & Clip for display; retain unclipped for downstream \\
\bottomrule
\end{tabular}
\end{center}

\section{Worked example}

5-hour session window: $t_{\text{start}} = 11{:}00$, $\Treset = 16{:}00$,
$\tnow = 13{:}00$, $\unow = 0.30$, so $\Dtrem = 3$\,h.

Last four snapshots ($\taurecent = 30$\,min):

\begin{center}
\begin{tabular}{@{}ccc@{}}
\toprule
$i$ & $t_i$ (h from 12{:}30) & $u_i$ \\
\midrule
1 & 0.000 & 0.270 \\
2 & 0.167 & 0.280 \\
3 & 0.333 & 0.295 \\
4 & 0.500 & 0.300 \\
\bottomrule
\end{tabular}
\end{center}

OLS: $\rols \approx 0.0630\,\mathrm{h}^{-1}$, $\seolssq \approx 6.30 \times
10^{-5}$.

Prior (suppose, from history, after the \S\ref{sec:rate} noise correction):
$\mu_0 = 0.080\,\mathrm{h}^{-1}$, $\tauzerosq = 3.6 \times 10^{-3}$.

Posterior: $\taupostsq \approx 6.19 \times 10^{-5}$, $\rpost \approx
0.0633\,\mathrm{h}^{-1}$ (data dominates).

Calibration: $\sigsessionsq = 2.5 \times 10^{-3}\,\mathrm{h}^{-1}$, $\bar\tau^2
= 3.6 \times 10^{-3}\,\mathrm{h}^{-2}$ (approximately matches $\tauzerosq$ in a
stable regime, as expected).

Forecast:
\[
  F = 0.30 + 0.0633 \cdot 3 \approx 0.490,
\]
\[
  \max(\taupostsq,\, \bar\tau^2) = \max(6.19 \times 10^{-5},\, 3.6 \times 10^{-3})
    = 3.6 \times 10^{-3} \quad \text{(floor active)},
\]
\[
  \sigma_F^2 \approx 9 \cdot 3.6 \times 10^{-3}
    + 3 \cdot 2.5 \times 10^{-3}
    \approx 3.24 \times 10^{-2} + 7.5 \times 10^{-3}
    \approx 4.0 \times 10^{-2},
\]
\[
  \sigma_F \approx 0.200.
\]

80\% CI: $[0.490 - 1.282 \cdot 0.200,\ 0.490 + 1.282 \cdot 0.200]
\approx [0.234,\ 0.746]$. Display: ``Projected 49\% by reset (80\% CI:
23\%--75\%)''.

Rate uncertainty dominates path noise by $4.3\times$. The recent slope is
locally well-estimated, but historically the rate has varied a lot across
sessions, and the v1.1 floor honors that.

ETA to 100\%: gap is 0.70, deterministic crossing at $13{:}00 +
0.70/0.0633 \approx 24{:}04$, far beyond reset. Returned as \texttt{nil}.
The MC confirms by producing infinite crossings for almost all
trajectories.

\section{Limitations and upgrade paths}

\begin{itemize}
  \item \textbf{Heavy-tailed pivots.} Real per-hour increments are
    heavier-tailed than Gaussian. The 80\% CI undercovers in the tails.
    Cheap fix when calibration data is plentiful: replace $z_{0.9}$ with an
    empirical quantile of standardized residuals.
  \item \textbf{Constant-rate assumption.} A mode shift within the window
    (deep work $\to$ meetings) breaks (A1). The recency window $\taurecent$
    absorbs slow drift but not abrupt shifts. The principled upgrade is a
    Kalman filter where $r$ is a slowly-varying latent process.
  \item \textbf{Global $\sigsessionsq$.} Pivot variance may depend on
    day-of-week or project type. Stratified calibration is the upgrade, but
    requires more sessions per stratum.
\end{itemize}

\end{document}
